\documentclass{reportkit}

\setreportkitleftheader{PRIVATE WEALTH MANAGEMENT CANDIDACY}
\setreportkitfooter{Analysis \& verification package | Pre-rewrite}
\setreportkitversion{v1.0}

\title{Private Wealth Management Candidacy\\Analysis \& Verification Package}
\author{Candidate: [Name redacted]}
\date{September 7, 2026}

\begin{document}
\maketitle

\begin{execsummary}
The candidate has a plausible path into a wirehouse or bank private wealth management (PWM) analyst program, but the current resume does not yet tell a coherent PWM story. The strongest potential differentiator is the combination of entrepreneurship and direct customer relationship work. That advantage is conditional on verifying the actual scope of the startup, particularly the claim that the candidate was involved in ``raising \$1.5 million in first-round funding.'' The largest structural weakness is the absence of direct PWM or financial-services experience; no wording change can substitute for an internship or part-time placement. The resume should therefore remain in pre-rewrite status until the startup facts are confirmed, while networking and internship outreach proceed in parallel.
\end{execsummary}

\begin{redflag}{Submission gate}
Do not finalize or submit the resume until the startup funding claim, client-acquisition scope, company status, and AI-system deployment status have been verified with the candidate. Every submitted claim should be defensible under direct interviewer follow-up.
\end{redflag}

\section{Purpose and decision scope}
This package separates three questions that should not be conflated: what the current resume already proves, what remains unknown and requires candidate confirmation, and how the eventual answers should change the resume strategy. It is a diligence document rather than a final resume rewrite.

\begin{researchproblem}{Primary question}
Does the candidate have a verified, personally owned entrepreneurship-and-client-relationship story strong enough to become the centerpiece of a PWM candidacy, or should the startup be de-emphasized in favor of other experience while direct financial-services experience is built?
\end{researchproblem}

\begin{assumption}{Target role and market}
The target is a wirehouse or bank Private Wealth Management Analyst Program, using Los Angeles as the high-cost U.S. market proxy. The assessment is grounded only in the resume facts supplied for this review; unverified claims are treated as unknown rather than inferred.
\end{assumption}

\section{Current-state assessment}

\subsection{What the resume currently demonstrates}

\begin{table}[htbp]
\centering
\caption{Current evidence and its relevance to PWM candidacy}
\begin{tabular}{p{0.22\textwidth} p{0.34\textwidth} p{0.34\textwidth}}
\toprule
\textbf{Dimension} & \textbf{Resume evidence} & \textbf{Assessment} \\
\midrule
Academic standing & 3.63 university GPA in Finance emphasis; 4.00 at prior community college; Dean's Honor List at both & Solid academic signal. It should not be a screening weakness, but it is not the core differentiator. \\
Client exposure & Startup references customer-acquisition meetings and managing new clients; grocery role references customer relations & Relevant PWM signal because relationship work matters, but current wording describes duties rather than scope, ownership, or outcomes. \\
Entrepreneurship & Undergraduate co-founder experience & Potentially unusual and differentiating if the commercial traction and personal contribution survive detailed questioning. \\
Direct PWM / financial services & None shown & Most important structural gap. Requires real experience rather than resume reframing. \\
Narrative coherence & JPMorgan IB Forage simulation, investment-banking framing, Wall Street Prep DCF modeling & Creates an IB-oriented signal inconsistent with the stated PWM target and may invite a ``fallback path'' question. \\
\bottomrule
\end{tabular}
\end{table}

\begin{evidencenote}{Interpretation}
The resume contains real relationship-oriented ingredients, but the strongest PWM narrative is latent rather than demonstrated. The key issue is not whether the document can be made to sound more client-oriented; it is whether the candidate actually sourced, pitched, closed, retained, or otherwise owned meaningful client and investor interactions.
\end{evidencenote}

\subsection{What cannot be assessed from the resume alone}
The current document does not establish whether the \$1.5 million figure was personally raised by the candidate, total company capitalization, pledged capital, a target, or another party's raise. It also does not establish the depth of customer acquisition, the company's current operating status, or whether the agentic AI system reached real deployment. These are factual gaps, not writing problems.

\begin{diagram}[
  type=evidence,
  caption={Evidence required before a strong startup-led PWM narrative is justified.},
  source={Conceptual diagram based on the supplied candidacy assessment, September 2026.},
  description={A three-tier evidence stack progressing from resume claims through verified ownership to measured commercial outcomes, showing that stronger positioning requires stronger evidence.}
]
\begin{evidencestack}[tiers=3]
  \evidencetier{Resume claim}
  \evidencetier{Verified personal ownership}
  \evidencetier[rk accent]{Measured commercial outcome}
\end{evidencestack}
\end{diagram}

\section{Verification protocol}
The questions below should be framed as normal pre-interview diligence. The objective is to establish precise, defensible facts before rewriting the experience section.

\subsection{Startup: highest priority}
\begin{enumerate}
  \item What did the company actually sell or provide to clients?
  \item How many paying clients did it have at peak?
  \item Did the company generate revenue? If so, what approximate amount or range can be disclosed?
  \item What was the candidate's specific role in customer-acquisition meetings: sourcing leads, pitching, negotiating, closing, maintaining relationships, or primarily attending meetings led by others?
  \item Walk through the ``\$1.5 million first-round funding'' claim:
  \begin{itemize}
    \item Who invested: angels, friends and family, venture capital, accelerator, or another source?
    \item Was the financing equity, debt, or another instrument?
    \item What did the candidate personally do to secure it: direct pitching, preparing materials, diligence support, or supporting another person's raise?
    \item Was \$1.5 million actually closed, or was it a target, projection, commitment, or other non-closed amount?
  \end{itemize}
  \item How many co-founders and employees did or does the company have?
  \item What is the company's current status: operating, paused, pivoted, or wound down? The resume reportedly shows an end date of May 2026; because this review is dated September 7, 2026, May 2026 is already past, so the issue to clarify is whether that date accurately reflects the end of the candidate's involvement and the company's present status.
  \item What was the ``Agentic AI system aimed at optimizing workflows''? Did it reach client or internal production use? Was there a measurable result such as time saved, throughput, or output volume?
\end{enumerate}

\begin{redflag}{Funding claim standard}
Treat ``raised \$1.5 million'' as a high-scrutiny claim. If the candidate did not personally play a material role in a completed \$1.5 million financing, the resume must use narrower language that accurately describes the candidate's contribution.
\end{redflag}

\subsection{Grocery retail role}
Identify specific examples of resolving a difficult customer situation, training a new employee, or solving an operational problem under pressure. These may be more valuable as behavioral-interview evidence than as additional resume bullets.

\subsection{Education and credentials}
Confirm the exact term and year of transfer from community college to university. Separately, determine whether the Securities Industry Essentials (SIE) exam has been scheduled or preparation has begun, and whether the candidate is willing to prioritize it over the next four to six weeks.

\begin{limitationnote}{Credential treatment}
Do not add the SIE to the resume until it has actually been passed. Preparation or intent is useful for planning, not as a completed credential.
\end{limitationnote}

\section{Decision framework after verification}

\begin{diagram}[
  type=flow,
  caption={Resume strategy after startup verification.},
  source={Conceptual diagram based on the supplied decision rules, September 2026.},
  description={A horizontal decision flow beginning with startup verification and branching to strong, softened, or peripheral evidence, each leading to a different resume treatment before a common PWM narrative review.}
]
\begin{reportflow}[direction=horizontal]
  \step{verify}{Verify startup facts}
  \step{strong}{Strong personal ownership}
  \step{soft}{Claim materially softer}
  \step{weak}{Role vague or peripheral}
  \step{lead}{Lead with startup}
  \step{narrow}{Narrow claim to facts}
  \step{deemph}{De-emphasize startup}
  \step{review}{Re-run PWM gap review}
  \branch{verify}{strong}
  \branch{verify}{soft}
  \branch{verify}{weak}
  \flowedge{strong}{lead}
  \flowedge{soft}{narrow}
  \flowedge{weak}{deemph}
  \merge{lead}{review}
  \merge{narrow}{review}
  \merge{deemph}{review}
\end{reportflow}
\end{diagram}

\subsection{Branch A: strong, specific, personally owned story}
If the candidate directly pitched clients, played a material role in a completed raise, and the company had real revenue or traction, rewrite the Experience section to lead with the startup. Entrepreneurship plus credible client-relationship ownership can become the centerpiece of the PWM candidacy. Build detailed STAR-format answers for ``tell me about your startup,'' ``why not investment banking,'' and ``why PWM'' from the same verified facts.

\subsection{Branch B: funding claim is materially softer than implied}
If \$1.5 million represents valuation, a target, pledged capital, another person's raise, or another materially different fact pattern, rewrite the bullet to describe only what the candidate personally did. For example, ``supported investor materials and meetings'' is preferable to an impressive but indefensible fundraising claim. Accuracy is the higher-value outcome because a single interviewer follow-up can otherwise undermine credibility across the entire resume.

\subsection{Branch C: role was vague, minimal, or peripheral}
If the candidate cannot demonstrate meaningful ownership, reduce the startup's prominence. Shift weight toward customer trust and operational reliability in the grocery role and toward leadership evidence. In this branch, securing relevant financial-services experience becomes a much larger lever than resume repositioning.

\begin{decisionpoint}{Default rule}
The strength of the resume claim must not exceed the strength of the verified evidence. Use the startup aggressively only when ownership and outcomes are both specific and defensible; otherwise narrow or de-emphasize it.
\end{decisionpoint}

\section{Actions that apply in every branch}
Trim or remove IB-signaling elements that weaken narrative coherence, particularly the JPMorgan IB Forage simulation entry if space is constrained or if it reinforces the impression that PWM is a fallback. Reframe relevant finance coursework and activities around markets, portfolio thinking, client communication, judgment, and relationship management only where the underlying facts support that framing.

Maintain explicit FLAG markers in working drafts for any metric that remains unconfirmed. Do not allow a flagged metric into a submission version. Continue internship and part-time placement outreach regardless of rewrite timing; direct PWM or adjacent financial-services experience remains the principal structural gap.

\begin{strategicpillars}[objective={Build a credible PWM candidacy}]
  \pillar{Verify}{Defensible claims and precise ownership}
  \pillar{Position}{Coherent client-and-relationship narrative}
  \pillar{Experience}{Direct PWM or adjacent placement}
  \pillar{Prepare}{SIE and interview stories}
\end{strategicpillars}

\section{Execution plan}

\begin{diagram}[
  type=roadmap,
  caption={Recommended execution sequence.},
  source={Conceptual diagram based on the supplied next steps, September 2026.},
  description={A three-horizon roadmap separating immediate fact verification, subsequent resume rewriting and gap analysis, and later finalization while networking proceeds in parallel.}
]
\begin{reportroadmap}[mode=horizons]
  \horizon{Now}{Verify}{Send diligence questions, Confirm startup facts, Confirm education dates}
  \horizon{Next}{Rewrite}{Apply decision branch, Rewrite Experience, Re-run gap analysis}
  \horizon{Then}{Submit}{Resolve all FLAGs, Finalize resume, Prepare interview stories}
\end{reportroadmap}
\end{diagram}

Networking and outreach should run concurrently with all three horizons. It does not depend on completion of the verification process.

\begin{deliverablenote}{Completion criteria}
The package is complete when the startup questions have direct answers, the applicable decision branch has been selected, all material metrics are verified or removed, the Experience section has been rewritten accordingly, and the revised resume has passed a fresh PWM gap analysis.
\end{deliverablenote}

\section{Risks and limitations}
This assessment is intentionally bounded by the supplied resume facts and does not independently verify the candidate's claims, recruiting-market conditions, firm-specific analyst-program requirements, or hiring timelines. It evaluates narrative credibility and evidence quality, not the probability of receiving an offer at a specific institution.

The absence of direct PWM experience is a structural constraint. A stronger resume can improve signal extraction from existing experience, but it cannot manufacture the industry exposure that many competing candidates may already possess.

\begin{limitationnote}{Date correction}
The source package described a May 2026 resume end date as being ``in the future relative to today.'' On the stated preparation date of September 7, 2026, May 2026 is in the past. This report corrects that inconsistency and retains the substantive diligence question: whether May 2026 accurately marks the candidate's end date and what the company's current status is.
\end{limitationnote}

\end{document}
